\documentclass[11pt,letterpaper]{article}
\usepackage[margin=0.85in]{geometry}
\usepackage{lmodern}
\usepackage{hyperref}
\hypersetup{colorlinks=false, urlcolor=black, pdfborder=0 0 0}
\setlength{\parindent}{0pt}
\setlength{\parskip}{0.45em}
\raggedright
\pagestyle{empty}

\begin{document}

\noindent
{\bfseries Ke-Cheng Hsu (Clio)}\\
Ph.D.\ Candidate, Department of Economics\\
The Hong Kong University of Science and Technology (HKUST)\\
(852) 6840\,9753 \quad $|$ \quad \href{mailto:khsu@connect.ust.hk}{khsu@connect.ust.hk}\\
\href{https://cliohsu1997.github.io/}{cliohsu1997.github.io} \quad $|$ \quad \href{https://github.com/cliohsu1997}{github.com/cliohsu1997}

\vspace{0.7em}
\noindent
October 26, 2026

\vspace{0.35em}
\noindent
Hiring Team\\
Analysis Group\\
Rue Montoyer 47\\
1000 Brussels\\
Re: Analyst Intern, Generalist, Brussels (2027 start; iCIMS 3045)

\vspace{0.2em}
Dear Hiring Team,

I am writing to apply for the Analyst Intern, Generalist, position in Analysis Group's Brussels office. I am a Ph.D.\ candidate in Economics at HKUST, working on empirical industrial organization and digital economics, and I expect to graduate in June 2027. I can start in July 2027. The length of the internship is flexible on my side.

I applied because of work like your support for Apple in Epic Games v.\ Apple. Epic claimed that Apple monopolized distribution and payment on the App Store. Your team, supporting Professor Richard Schmalensee, treated the App Store as a two-sided platform and asked which side the rules help and which side pays. Platform rules and who bears the cost are the questions I study.

My job market paper asks a version of that question for an online seller: should a Taobao apparel seller subsidize the return shipping fee for every buyer, or let buyers pay for that protection themselves? With return rates around 20--30\%, the answer moves real money. Using Taobao apparel transactions (2017--2023; six apparel categories), I clean large transaction data, estimate a structural model of purchase, return-shipping-insurance take-up, and return decisions, and rank policies by seller profit, consumer surplus, and social welfare. With Professors Kohei Kawaguchi and Kei Kawai, I study pricing and inventory under resale price maintenance, combining descriptive evidence, regressions, and a dynamic oligopoly model built from monthly bookstore point-of-sale records across 1{,}000+ stores.

I work in R, Stata, Python, and C++ on AWS and HPC/SLURM, and I build pipelines I can rerun from raw data to final numbers, including figures. I have also used Python to scrape product pages and OCR product images. As the lecturer for HKUST's MSc Math Camp (4.53/5) and a TA for PhD and undergraduate economics (4.72/5 on average), I am used to turning a technical result into something a non-specialist can use.

My CV is attached. I can also send an unofficial transcript and my job market paper. Thank you for reading; I hope we get the chance to talk.

\vspace{0.7em}
Sincerely,\\[1.0em]
Ke-Cheng Hsu (Clio)

\end{document}
